% 258
% \S8 МНОГОГРАННИКИ Понятие многогранника, Элементы многогранника

\par{1. Определение многогранника.}
Многогранник --- ограниченное тело, поверхность которого состоит из конечного
числа многоугольников
\footnote{Здесь многоугольником мы называем конечную часть плоскости,
ограниченную одной или несколькими замкнутыми ломанными.}.

На рисунке \ref{fig:8_1_1} изображено несколько многогранников.

\begin{figure}\label{fig:8_1_1}
% стр 258 рис 1 строенный
\caption{Параллелепипед}
\caption{Призма}
\caption{Пирамида}
\caption{Параллелепипед с~прорезанным квадратным отверстием}
\caption{Много составленный из трёх кубов}
\end{figure}

% 259
\par{2. Элементы многогранника.}
Многоугольники, из которых составлена поверхность многогранника называются
его гранями. При этом граница грани может состоять из нескольких замкнутых
ломаных. На рисунках \ref{fig:8_1_1_г} и~\ref{fig:8_1_1_д} такого типа грани
указаны стрелками.
\begin{note}Замечание.
Мы сказали, что грани многогранника --- это многоугольники, из которых составлена
поверхность многогранника. Это понятие следует уточнить следующим образом:
многоугольник на поверхности многогранника называется его гранью, если он
не является частью другого многоугольника, так же лежащего на поверхности
многогранника.
\end{note}

% 260
Без этого уточнения можно было бы, например, называть гранью куба любой
многоугольник, расположенный на его поверхности так, как это изображено
на рисунке \ref{fig:8_1_2} (этот многоугольник заштрихован.

\begin{figure}\label{fig:8_1_2}
% стр 260 рис 2
\caption{Рис.2. Многоугольник $PQRST$ расположен на поверхности куба, но не
является его гранью, так как является частью многоугольника $ABCD$ (квадрата),
также расположенного на: поверхности куба. Очевидно, квадрат $ABCD$ --- это
грань куба.}
\end{figure}

Стороны граней называются ребрами многогранника, а вершины рёбер с~вершинами
многогранника.

Грани, ребра и~вершины многогранника называются его элементами.

На рисунке \ref{fig:8_1_3} изображен тетраэдр и~выделены его элементы.

\begin{figure}\label{fig:8_1_3}
% стр 260 рис 3
\caption{Рис.З. Элементы тетраэдра: видимые грани заштрихованы, видимые ребра
изображены сплошными отрезками, невидимое ребро --- штриховым,
вершины изображены кружочками.}
\end{figure}

% 261
\par{3 Многогранная поверхность. Развертки многогранника.}

Многогранной поверхностью называется поверхность, составленная из многоугольников
(рис.\ \ref{fig:8_1_4}).

\begin{figure}\label{fig:8_1_4}
% стр 261 рис 4
\caption{Рис.4. Многогранная поверхность. Она составлена из многоугольников.}
\end{figure}

Примером замкнутой многогранной поверхности может служить поверхность
многогранника.

Процесс составления многогранной поверхности из многоугольников можно представить
себе как операцию склеивания этой поверхности из многоугольников.
На рисунке \ref{fig:8_1_5} показано, как из четырёх треугольников склеивается
(составляется) многогранная поверхность --- тетраэдр. Стрелками указано,
какие стороны треугольников склеивают друг с другом.

\begin{figure}\label{fig:8_1_5}
% стр 261 рис 5
\caption{Рис.5 Поверхность правильного тетраэдра может быть
склеена, из четырех одинаковых равносторонних треугольников. Стрелками
указано, какие стороны этих треугольников склеиваются друг с другом.}
\end{figure}

% 262

\begin{figure}\label{fig:8_1_x}
% стр 262 рис без номера
\caption{Правильный тетраэдр составлен (склеен) из четырех равносторонних
треугольников.}
\end{figure}

Набор многоугольников, из которых склеивается многогранная поверхность, называется
разверткой этой поверхности (развёрткой многогранника). На рисунках
\ref{fig:8_1_6} и~\ref{fig:8_1_7} указаны развертки куба и~правильной пятиугольной
пирамиды. При этом стороны) многоугольников, которые склеиваются друг с~другом,
обозначены одинаковыми буквами.

\begin{figure}\label{fig:8_1_6}
% стр 262 рис 6
\caption{Развёртка куба и~куб, составленный (склеенный) из этой развёртки.
Одинаковыми буквами обозначены стороны многоугольников развёртки, которые
склеиваются друг с~другом.}
\end{figure}

% 263

\begin{figure}\label{fig:8_1_7}
% стр 263 рис 7
\caption{Рис.7. Развертка правильной пятиугольной пирамиды. Склеиваемые друг
с~другом стороны обозначены одинаковыми буквами.}
\end{figure}

Иногда в~многоугольники развертки данного многогранника склеивают так, что
в~результате получают его  развёртку в~виде одного многоугольника. При этом,
конечно, указывается, какие свободные (не склеенные) стороны этого многоугольника
должны быть склеены друг с~другом. На рисунке \ref{fig:8_1_8} представлена
такая развертка куба, полученная из развертки, приведённой на`рисунке \ref{fig:8_1_6},
посредством попарного` склеивания сторон квадратов, обозначенных буквами
$b$, $d$, $h$, $k$, $l$. При этом склеенные стороны выделены на рисунке
\ref{fig:8_1_8} штриховыми линиями.

\begin{figure}\label{fig:8_1_8}
% стр 263 рис 8
\caption{Рис. 8. Развертка куба в виде одного многоугольника.}
\end{figure}

% 264

На рисунке \ref{fig:8_1_9} изображена развертка правильной пятиугольной пирамиды.
Эта развертка получена из развертки, данной на рисунке \ref{fig:8_1_7},
путём попарного склеивания по сторонам многогранников обозначенных буквами
$a$, $b$, $c$, $d$ и~$e$. Склеенные стороны выделены на рисунке \ref{fig:8_1_9}
штриховыми линиями.

\begin{figure}\label{fig:8_1_9}
\caption{Рис. 9.Развертка правильной пятиугольной пирамиды в виде одного
многоугольника.}
\end{figure}

\begin{note}
Замечание. У~заданного многоугольника имеется бесконечно много разверток.
Каждая из его разверток может быть получена из другой, образно говоря, путём
перекройки --- одну из разверток можно разрезать на треугольники так, что путём
склеивания по определённым сторонам полученных треугольников можно перейти
к~заданной развертке. Доказать это несложно. Попробуйте это сделать самостоятельно.
Достаточно ребра одной развёртки окрасить, например, в~красный цвет, а~ребра
другой --- в~синий, и~обе развертки уложить на многогранник.
Ну, а что нужно сделать дальше --- подумайте сами.
\end{note}
